% ============================================================
%  LÍNEA DE AIRE DE ENTRADA — secuencia 1 a 6 (esquemático)
%  Sin librería calc: solo coordenadas explícitas.
% ============================================================
\begin{tikzpicture}[font=\scriptsize, line cap=round, line join=round]
\colorlet{aire}{red!80!black}

% ---------- mangueras (debajo de todo) ----------
\draw[aire, line width=1.4pt] (1.9,2.0) -- (4.3,2.0);            % bomba -> filtro
\draw[aire, line width=1.4pt] (5.7,2.0) -- (8.1,2.0);            % filtro -> PTC
\draw[aire, line width=1.4pt] (9.5,2.0) -- (10.0,2.0) -- (10.0,2.9)
      -- (11.9,2.9) -- (11.9,2.6);                               % PTC -> humidificador (por arriba)
\draw[aire, line width=1.4pt] (12.8,2.2) -- (13.3,2.2) -- (13.3,1.75)
      -- (14.55,1.75) -- (14.55,1.55);                           % humidificador -> rotámetro (entra abajo)
\draw[aire, line width=1.4pt] (14.8,3.3) -- (15.4,3.3)
      -- (15.4,2.0) -- (16.6,2.0);                               % rotámetro -> válvula
\draw[aire, line width=1.4pt, -{Stealth[length=3mm]}]
      (17.3,2.0) -- (18.2,2.0) -- (18.2,1.65) -- (18.7,1.65);    % válvula -> plenum

% ---------- aire ambiente ----------
\draw[aire, line width=1.4pt, -{Stealth[length=3mm]}] (-0.6,2.0) -- (0.5,2.0);
\node[anchor=east, text=aire] at (-0.7,2.0) {aire ambiente};

% ---------- 1. bomba ----------
\draw[fill=blue!12] (0.5,1.6) rectangle (1.9,2.4);
\draw[fill=blue!35] (1.2,2.0) circle (0.22);
\node[circle, fill=red, text=white, font=\scriptsize\bfseries, inner sep=1.5pt] at (1.2,3.0) {1};
\node[anchor=north, align=center] at (1.2,1.45) {bomba de aire 10 L/min\\(siempre encendida)};

% ---------- 2. filtro HEPA ----------
\draw[fill=gray!20] (4.3,1.55) rectangle (5.7,2.45);
\foreach \x in {4.7,5.0,5.3} \draw[gray!70] (\x,1.62) -- (\x,2.38);
\node[circle, fill=red, text=white, font=\scriptsize\bfseries, inner sep=1.5pt] at (5.0,3.05) {2};
\draw[gray!60] (5.0,1.55) -- (5.0,1.0);
\node[anchor=north, align=center] at (5.0,0.95) {filtro HEPA\\(polvo y esporas fuera)};

% ---------- 3. resistencia PTC ----------
\draw[fill=orange!25] (8.1,1.7) rectangle (9.5,2.3);
\draw[orange!80!black, line width=1.1pt] (8.25,1.9) -- (8.45,2.1) -- (8.65,1.9)
      -- (8.85,2.1) -- (9.05,1.9) -- (9.25,2.1) -- (9.4,2.0);
\node[circle, fill=red, text=white, font=\scriptsize\bfseries, inner sep=1.5pt] at (8.8,2.95) {3};
\node[anchor=north, align=center] at (8.8,1.45) {resistencia PTC 12 V\\(calienta si hace falta)};

% ---------- 4. humidificador de burbuja ----------
\draw[fill=cyan!8] (11.4,1.35) rectangle (12.8,2.6);
\fill[cyan!25] (11.46,1.41) rectangle (12.74,2.05);
\draw[cyan!60!black] (11.46,2.05) -- (12.74,2.05);
\draw[aire, line width=1.2pt] (11.9,2.6) -- (11.9,1.5);
\foreach \y/\x in {1.6/12.0,1.75/12.12,1.92/11.98} \draw[cyan!60!black] (\x,\y) circle (0.045);
\node[circle, fill=red, text=white, font=\scriptsize\bfseries, inner sep=1.5pt] at (12.1,3.35) {4};
\draw[gray!60] (12.1,1.35) -- (12.1,1.0);
\node[anchor=north, align=center] at (12.1,0.95) {humidificador de burbuja\\(el aire entra húmedo)};

% ---------- 5. rotámetro (vertical) ----------
\draw[fill=cyan!10] (14.3,1.55) rectangle (14.8,3.65);
\draw[-{Stealth[length=2mm]}, gray!70] (14.55,1.7) -- (14.55,3.4);
\fill[red!70] (14.41,2.45) rectangle (14.69,2.65);
\foreach \y in {2.0,2.5,3.0} \draw (14.8,\y) -- (14.95,\y);
\node[circle, fill=red, text=white, font=\scriptsize\bfseries, inner sep=1.5pt] at (14.55,4.05) {5};
\node[anchor=north, align=center] at (14.55,1.45) {rotámetro (lectura de Q)\\va vertical};

% ---------- 6. válvula solenoide ----------
\draw[fill=gray!25] (16.6,1.8) -- (16.6,2.2) -- (16.95,2.0) -- cycle;
\draw[fill=gray!25] (17.3,1.8) -- (17.3,2.2) -- (16.95,2.0) -- cycle;
\draw[fill=white] (16.7,2.25) rectangle (17.2,2.7);
\node[font=\tiny] at (16.95,2.48) {12 V};
\node[circle, fill=red, text=white, font=\scriptsize\bfseries, inner sep=1.5pt] at (16.95,3.15) {6};
\draw[gray!60] (16.95,1.8) -- (16.95,1.0);
\node[anchor=north, align=center] at (16.95,0.95) {válvula solenoide\\(modula la aireación)};

% ---------- panera (destino) ----------
\draw[fill=orange!6] (18.7,1.45) rectangle (20.3,2.65);
\fill[gray!15] (18.74,1.49) rectangle (20.26,1.85);
\node[font=\tiny, text=gray!60] at (19.5,1.67) {PLENUM};
\draw[fill=orange!40] (18.85,2.0) rectangle (19.45,2.4);
\draw[fill=orange!40] (19.55,2.0) rectangle (20.15,2.4);
\node[anchor=north, align=center] at (19.5,1.3) {al plenum\\de la panera};

\end{tikzpicture}
